\documentclass{thesisclass}

%% -------------------------------
%% |    IM Thesis Template       |
%% -------------------------------
%% Further additions by: Philipp Stroehle, IM, 2013 
%% philipp.stroehle "at" kit.edu

%% Notes:
%% Language switch after \begin{document}

% Based on thesisclass.cls of Timo Rohrberg, 2009
% ----------------------------------------------------------------
% Thesis - Main document
% ----------------------------------------------------------------

% Select your thesis language here:
\usepackage[english]{babel}
%\usepackage[ngerman]{babel}


%% ---------------------------------
%% |      Additional packages      |
%% ---------------------------------
%% 

\usepackage{graphicx}
%http://en.wikibooks.org/wiki/LaTeX/Importing_Graphics#Graphics_storage
\DeclareGraphicsExtensions{.pdf,.png,.jpg}
\graphicspath{{./figures/}} %Use curly braces for each path to add and don't
% forget trailing slash '/'
% \usepackage{epstopdf} %Nice to automatically convert eps figures to pdf
% format  (from inkscape, etc)
\usepackage[square,sort,comma,numbers]{natbib}
\usepackage{booktabs}
\usepackage{amsmath}
\usepackage{amsthm}
\usepackage{amsfonts}
\usepackage{amssymb}
\usepackage{mathtools}
\usepackage{tikz}
\usepackage{booktabs}
\usepackage{multirow}
\usepackage{float}
\usetikzlibrary{positioning}
\usetikzlibrary{calc}
\usepackage[algo2e,linesnumbered,ruled]{algorithm2e}
\usepackage{pdflscape} 
\usepackage{placeins}
\usepackage[procnames]{listings}
\usepackage{color}
\usepackage[draft]{todonotes}
\usepackage{csquotes}
	
\usetikzlibrary{positioning}
\tikzset{main node/.style={circle,fill=blue!20,draw,minimum size=1cm,inner sep=0pt},}

\usetikzlibrary{arrows}


%% ---------------------------------
%% | Needed for the List of Abbreviations |
%% ---------------------------------
\usepackage{nomencl}
\renewcommand{\nomname}{List of Abbreviations}
% Punkte zw. Abkürzung und Erklärung
\setlength{\nomlabelwidth}{.32\hsize}
%.32\hsize bezieht sich auf die Punkte
% zwischen Abkürzung und Erklärung -> je größer die Zahl, desto mehr
% Punkte
\renewcommand{\nomlabel}[1]{#1 \dotfill}
% Zeilenabstäde verkleinern
\setlength{\nomitemsep}{-\parsep}
\makenomenclature
\usepackage[normalem]{ulem}
\newcommand{\markup}[1]{\uline{#1}}

%% ---------------------------------
%% |      Theorems and definitions |
%% ---------------------------------
%% 

\newtheorem{definition}{Definition}[chapter]
\newtheorem{proposition}{Proposition}[chapter]
%% ---------------------------------
%% | Information about the thesis  |
%% ---------------------------------

% uncomment one of the following, according to your thesis
\newcommand{\mytype}{\iflanguage{english}{Master's Thesis}{Masterarbeit}} 
%\newcommand{\mytype}{\iflanguage{english}{Master's Thesis}{Masterarbeit}} 
%\newcommand{\mytype}{\iflanguage{english}{Seminar Thesis}{Seminararbeit}} 

\newcommand{\myname}{Patrick Bisenius}
\newcommand{\matricle}{1640015}
\newcommand{\mytitle}{\iflanguage{english}{Computing Top-k
Closeness Centrality in Fully-dynamic Graphs}{Berechnung von
Top-k-Nähezentralität in volldynamischen Graphen}}
\newcommand{\myinstitute}{\iflanguage{english}
{At the Department of Informatics\\Institute of Theoretical Informatics}
{Institut für Theoretische Informatik (ITI)}
}

\newcommand{\reviewerone}{Jun. Prof. Dr. Henning Meyerhenke}
\newcommand{\reviewertwo}{?}
\newcommand{\advisor}{Elisabetta Bergamini}
\newcommand{\advisortwo}{}

\newcommand{\timestart}{\iflanguage{english}{1st December 2016}{1. Januar
2016}}
\newcommand{\timeend}{\iflanguage{english}{31st May 2017}{31. Mai 2017}}
\newcommand{\submissiontime}{02. Mai 2017}

%% -------------------------------
%% |  Information for PDF file   |
%% -------------------------------
%% IM: Auto-Fill this information
\hypersetup{
 pdfauthor={\myname},
 pdftitle={\mytitle},
 pdfsubject={\mytype},
 pdfkeywords={\mytype}
}

%% ---------------------------------
%% | ToDo Marker - only for draft! |
%% ---------------------------------
% Remove this section for final version!
\setlength{\marginparwidth}{20mm}

%\newcommand{\margtodo}
%{\marginpar{\textbf{\textcolor{red}{ToDo}}}{}}

%\newcommand{\todo}[1]
%{{\textbf{\textcolor{red}{(\margtodo{}#1)}}}{}}


%% --------------------------------
%% | Old Marker - only for draft! |
%% --------------------------------
% Remove this section for final version!
\newenvironment{deprecated}
{\begin{color}{gray}}
{\end{color}}


%% --------------------------------
%% | Settings for word separation |
%% --------------------------------
% Help for separation:
% In german package the following hints are additionally available:
% "- = Additional separation
% "| = Suppress ligation and possible separation (e.g. Schaf"|fell)
% "~ = Hyphenation without separation (e.g. bergauf und "~ab)
% "= = Hyphenation with separation before and after
% "" = Separation without a hyphenation (e.g. und/""oder)

% Describe separation hints here:
\hyphenation{
% Pro-to-koll-in-stan-zen
% Ma-na-ge-ment  Netz-werk-ele-men-ten
% Netz-werk Netz-werk-re-ser-vie-rung
% Netz-werk-adap-ter Fein-ju-stier-ung
% Da-ten-strom-spe-zi-fi-ka-tion Pa-ket-rumpf
% Kon-troll-in-stanz
Lauf-zeit-ver-bes-se-run-gen
}


%% ------------------------
%% |    Including files   |
%% ------------------------
% Only files listed here will be included!
% Userful command for partially translating the document (for bug-fixing e.g.)
\includeonly{%
titlepage,
abstract,
text/abbreviations,
text/introduction,
text/content,
text/evaluation,
text/conclusion,
text/declaration,
text/appendix
}

%%%%%%%%%%%%%%%%%%%%%%%%%%%%%%%%%
%% Here, main documents begins %%
%%%%%%%%%%%%%%%%%%%%%%%%%%%%%%%%%
\begin{document}

\frontmatter
\pagenumbering{roman}
%% titlepage.tex
%%

% coordinates for the bg shape on the titlepage
\newcommand{\diameter}{20}
\newcommand{\xone}{-15}
\newcommand{\xtwo}{160}
\newcommand{\yone}{15}
\newcommand{\ytwo}{-253}

\begin{titlepage}
% bg shape
\begin{tikzpicture}[overlay]
\draw[color=gray]  
 		 (\xone mm, \yone mm)
  -- (\xtwo mm, \yone mm)
 arc (90:0:\diameter pt) 
  -- (\xtwo mm + \diameter pt , \ytwo mm) 
	-- (\xone mm + \diameter pt , \ytwo mm)
 arc (270:180:\diameter pt)
	-- (\xone mm, \yone mm);
\end{tikzpicture}
	\begin{textblock}{10}[0,0](4,2.5)
		\iflanguage{english}
			{\includegraphics[width=.3\textwidth]{logos/KITLogo_EN_RGB.pdf}}
			{\includegraphics[width=.3\textwidth]{logos/KITLogo_DE_RGB.pdf}}		
	\end{textblock}
	\begin{textblock}{10}[0,0](10,2.5)
		%\includegraphics[width=.6\textwidth]{logos/IISMLogo_RGB.pdf}
	\end{textblock}
	\changefont{ppl}{m}{n}	% helvetica	(phv), % IM Style: palatino (ppl) 
	\vspace*{3.5cm}
	\begin{center}
		\Huge{\mytitle}
		\vspace*{2cm}\\
		\Large{
			\iflanguage{english}{\mytype}			
								{\mytype}
		}\\
		\vspace*{1cm}
		\huge{\myname}\\
			\Large{\matricle}\\
		\vspace*{1cm}
		\Large{		
			\myinstitute
		}
	\end{center}
	\vspace*{3cm}
\Large{
\begin{center}
\begin{tabular}[ht]{l c l}
  % Gutachter sind die Professoren, die die Arbeit bewerten. 
  %\iflanguage{english}{Reviewer}{Erstgutachter}: & \hfill  & \reviewerone\\
  \iflanguage{english}{Reviewer}{Gutachter}: & \hfill  & \reviewerone\\
  %\iflanguage{english}{Second reviewer}{Zweitgutachter}: & \hfill  & \reviewertwo\\
  \iflanguage{english}{Advisor}{Betreuender Assistent}: & \hfill  & \advisor\\
  % \iflanguage{english}{Second advisor}{Zweiter betreuender Mitarbeiter}: & \hfill  & \advisortwo\\
  % IM: No second advisor
  % Der zweite betreuende Mitarbeiter kann weggelassen werden. 
\end{tabular}
\end{center}
}


\vspace{2cm}
\begin{center}
\large{\timeend}
% \iflanguage{english}{Duration:}{Bearbeitungszeit:} \timestart \hspace*{0.25cm}
% -- \hspace*{0.25cm} 
% \timeend}
\end{center}


\begin{textblock}{10}[0,0](4,16.8)
\tiny{ 
	\iflanguage{english}
		{KIT -- University of the State of Baden-Wuerttemberg and National Laboratory of the Helmholtz Association}
		{KIT -- Universit\"at des Landes Baden-W\"urttemberg und nationales Forschungszentrum der Helmholtz-Gesellschaft}
}
\end{textblock}

\begin{textblock}{10}[0,0](14,16.75)
\large{
	\textbf{www.kit.edu} 
}
\end{textblock}

\end{titlepage}




\chapter{Abstract}

Closeness is a widely-studied centrality measure in network analysis. Since it requires all pairwise distances, computing closeness centrality for all nodes is unfeasible for large real-world networks. However, for most applications, it is only necessary to find the $k$ most central nodes and not all closeness values. Prior work has shown that computing the top-k nodes with highest closeness can be done much faster than computing closeness for all nodes in real-world networks.

However, for networks that evolve over time, no dynamic top-$k$ closeness algorithm exists that improves on static recomputation. In this thesis, we present several techniques that allow us to efficiently recompute the $k$ nodes with highest closeness after an edge insertion or an edge deletion. Our algorithms use information obtained during earlier computations to skip unnecessary work. However, they do not require asymptotically more memory than the static algorithms (i.e. linear in the number of nodes). We propose separate algorithms for complex networks and networks with large diameter, such as street networks. For edge insertions in undirected complex networks, our average speedup on recomputation is about 40 and it is up to 160 for some graphs. For street networks, the average speedup is 70 and it is up to 150. For some edges, our algorithm is up to three orders of magnitude faster than static recomputation.

We also study the problem of finding the group of $k$ nodes with high closeness. We adapt a static greedy algorithm for the dynamic case. We use our dynamic algorithms for top-k closeness centrality and employ similar techniques to reduce the amount of computation. On average, our dynamic algorithm is between 8 and 15 times faster, depending on the group size.
\vspace{1cm}


\chapter{Zusammenfassung}\label{chap:abstract_german}

Nähezentralität ist eine weit verbreitete Metrik im Feld der Netzwerkanalyse. Da zur Berechnung die Distanzen aller Knotenpaare nötig wären, ist es unpraktikabel, die Nähezentralität aller Knoten in großen realen Netzwerken zu berechnen. Für einige Anwendungen reicht es allerdings aus, die $k$ zentralsten Knoten zu bestimmen und nicht alle Nähezentralitäten. Frühere Arbeiten haben bereits gezeigt, dass die Berechnung der $k$ zentralsten Knoten wesentlich schneller durchgeführt werden kann als die Berechnung der Nähezentralität für alle Knoten. 

Allerdings existieren bisher keine Algorithmen für Netzwerke, die sich im Laufe der Zeit verändern, die besser als eine statische Neuberechnung sind. Im Rahmen dieser Arbeit stellen wir nun mehrere Techniken vor, die uns erlauben, die $k$ zentralsten Knoten eines Netzwerks nach Einfügen oder Löschen einer Kante auf effiziente Weise neuzuberechnen. Unsere Algorithmen verwenden dabei Information aus früheren Berechnungen um unnötigen Rechenaufwand zu vermeiden. Dafür benötigen unsere Algorithmen nur linearen zusätzlichen Speicher, also asymptotisch gesehen nicht mehr als die statischen Algorithmen. Wir stellen verschiedene Algorithmen für komplexe Netzwerke und Netzwerke mit großen Durchmesser, wie zum Beispiel Straßennetzwerke vor. Unser dynamischer Algorithmus ist im Schnitt etwa 40 mal schneller als die statische Referenz für ungerichtete komplexe Netzwerke, in einzelnen Fällen bis zu 160 mal schneller. Für Straßennetzwerke wird die Berechnung im Schnitt um den Faktor 70 beschleunigt, für einzelne Graphen sogar um einen Faktor von 150. Für einige bestimmte Kanten ist unser dynamischer Algorithmus um bis zu drei Größenordnungen schneller als eine statische Neuberechnung.

Wir beschäftigen uns außerdem mit dem Problem der Gruppenzentralität. Dabei geht es darum, eine Gruppe mit $k$ Knoten zu finden, die als Ganzes eine hohe Zentralität hat. Dabei passen wir einen statischen Greedy-Algorithmus für dynamische Graphen an. Wir verwenden dabei unter anderem unseren dynamischen Algorithmus für Top-$k$-Nähezentralität und adaptieren weitere ähnliche Techniken um den Rechenaufwand zu reduzieren. Im Durchschnitt ist unser dynamischer Algorithmus acht bis fünfzehn Mal schneller als die statische Referenz, je nachdem wie groß die gesuchte Gruppe ist.

\chapter{Acknowledgements}

First of all, I would like to thank Professor Meyerhenke for giving me the opportunity to work on this very interesting topic.

My special thanks goes to my advisor Elisabetta Bergamini. She helped me find the topic and assisted me in every way I needed. We had many productive discussions that produced the ideas and solutions presented in this thesis.

I would like to thank all of my friends who were very accepting when I did not have the time or mood to attend social events.

At last, I would like to thank my parents, my brother, and my sister for supporting me all this time.
% IM Style: No additional blank page
% \blankpage


%% -------------------
%% |   Directories   |
%% -------------------
\tableofcontents
% IM Style: No additional blank page
% \blankpaage

% Do not include a list of figures, list of tables and list of abbreviations,if the work is a seminar thesis
\iflanguage{english}{
\ifthenelse{\equal{\mytype}{Seminar Thesis}}
{}
{
\listoffigures \addcontentsline{toc}{chapter}{List of Figures} 
\listoftables  \addcontentsline{toc}{chapter}{List of Tables} 
%\printnomenclature   \addcontentsline{toc}{chapter}{List of Abbreviations} 
}
}
{
\ifthenelse{\equal{\mytype}{Seminararbeit}}
{}
{
\listoffigures \addcontentsline{toc}{chapter}{Abbildungsverzeichnis} 
\listoftables \addcontentsline{toc}{chapter}{Tabellenverzeichnis} 
%\printnomenclature \addcontentsline{toc}{chapter}{Abkürzungsverzeichnis}
}
}

% Execute this command for index creation, i.e., for abbreviations by the nomencl package

%makeindex thesis.nlo -s nomencl.ist -o thesis.nls



%\printnomenclature



%% -----------------
%% |   Main part   |
%% -----------------
\mainmatter
\pagenumbering{arabic}
%% introduction.tex
%%

%% ==============================
\chapter{Introduction}
\label{ch:Introduction}
%% ==============================

The concept of a \emph{network} is used as a tool to model interactions and connections in several fields of science. ``The scientific study of networks \textelp{} is an interdisciplinary field that combines ideas from mathematics, physics, biology, computer science, the social sciences and many other areas'', Mark Newman writes in his book \emph{Networks. An introduction}~\cite{newmannetworks}.   

One fundamental concept in network analysis is \emph{Centrality}. Centrality measures are used to ascertain the importance of nodes within a graph. Practical applications include finding the most influential people in a social network~\cite{FREEMAN1978215}, identifying key infrastructure nodes in computer networks, analyzing the effects of human land use to organism movement~\cite{estrada2008using}, or finding so-called super-spreaders in disease transmission networks~\cite{dekker2013network}. 

One of the most widely-known centrality measures is \emph{PageRank}, proposed by the founders of Google, Sergey Brin and Larry Page~\cite{page1999pagerank}. The idea behind PageRank is that a webpage is important if it is linked to by other important webpages. Instead of simply relying on the absolute number of links to a webpage (a metric that was often used in citation networks), the quality of the links is taken into account. PageRank is similar to Eigenvector centrality. At the time, PageRank was a considerable improvement over other search engines.

Another popular centrality measure is called Betweenness Centrality. A node is important if it is part of many shortest paths between other node pairs. For instance, this can be useful to find the routers in a computer network that are most integral to its stability.
\paragraph{Closeness centrality}

Closeness centrality is based on the intuition, first presented by Alexander Bavelas in 1950~\cite{bavelas1950communication}, that a node is important if its distance to other nodes in the graph is small. There are different definitions of closeness centrality that are applicable in different contexts. For strongly connected graphs, one can simply compute the sum of the distances from one node to all the other nodes. However, this approach leads to problems on disconnected graphs. It is not initially clear how to treat nodes which are not connected by a path in the graph. Simply assuming an infinite (or arbitrarily large) distance would completely distort the resulting closeness values of affected nodes. This is because closeness centrality involves computing the sum of distances to other nodes. Arbitrarily large or infinite terms automatically lead to a sum total that is arbitrarily large and therefore loses any informative value. One possible solution for this problem is to only take into account the distances of reachable nodes, then scaling the result with the number of reachable nodes. Another approach is called \emph{harmonic closeness centrality} which has shown to have nice properties on arbitrary graphs~\cite{boldi2014axioms}. Instead of summing all distances and then computing the inverse, the harmonic closeness centrality is obtained by computing the sum of the inverse distances between nodes. Disconnected node pairs do not contribute to the total sum of inverse distances.

Computing the closeness centrality of a node in an unweighted graph requires a complete breadth-first search (BFS), and a complete run of Dijkstra's algorithm~\cite{dijkstra1959note} with weighted graphs. It requires solving the \emph{all-pairs-shortest-path} problem to compute the closeness centrality of each node in the graph. The computational effort for this is often impractical, especially on large real-world networks. Moreover, this effort can currently not be avoided if the application requires an exact ranking of all nodes by their exact closeness centrality. \todo{Cite survey of different centrality measures} 

For some applications, however, it is enough to compute a list of the $k$ most central nodes. This problem is called \emph{Top-k closeness}. Limiting the problem to the $k$ most central nodes can decrease the required computational effort significantly. It is only necessary to compute the exact closeness centrality of the $k$ most central nodes to rank them. For all other nodes, it is sufficient to obtain an upper bound for their closeness that is smaller than the exact closeness of the $k$-th most central node.

%\todo{Mention algorithms for dynamic closeness}

\paragraph{Dynamic Top-k closeness centrality}

In some cases, it is enough to compute closeness centralities only once because the underlying graph is static. Now consider a social network which constantly adds new users, which is effectively a node insertion in the underlying graph, and existing users befriend other users, which is an edge insertion. Terminating a friendship in the social network corresponds to an edge removal.

Each modification of an undirected graph affects at least the closeness centralities of the directly affected nodes, that is, the nodes incident to a newly inserted or removed edge. In directed graphs, only the target node of the directed edge is directly affected. It is also possible that there are new shortest paths between pairs of nodes that use the newly inserted edge. Analogously, removing an edge might increase the distance between node pairs because every shortest path between them contained the removed edge. A simple strategy to get the new closeness centralities of each node is to re-run the static algorithm on the modified graph, ignoring any information collected by previous runs of the algorithm.

\paragraph{Group closeness}
The concept of closeness centralities for single nodes can be extended to groups of nodes. The distance between a node $v$ and a group $S$ is defined as the smallest distance between $v$ and any node of the group. The problem to find a group of size $k$ such that the total distance of all nodes in the graph to the group is minimal is called the \emph{maximum closeness centrality group identification} (MCGI) problem by Chen et al. in~\cite{Chen2016}. Since the problem is shown to be NP-hard, no efficient exact algorithm exists at this point. However, there are approximative greedy algorithms~\cite{Chen2016,zhao2014measuring} for the problem. Bergamini and Meyerhenke improve on the work in~\cite{Chen2016} by reducing the memory requirements and total number of operations \todo{Cite unpublished work}.


\paragraph{Contributions}
This thesis contributes a dynamic algorithm for Top-k closeness that handles both edge insertions and edge removals. It is based on the static algorithm first proposed by Borassi et al for complex networks in ~\cite{borassi2015fast}, and the additional optimizations for street networks proposed by Bergamini et al. in ~\cite{bergamini2016computing}. Our algorithm re-uses information obtained by an initial run of the static algorithm and tries to skip the re-computation of closeness centralities for nodes that are unaffected by modifications of the graph. In some cases, even the upper bounds for the closeness centralities of affected nodes can be updated with little computational effort. Our dynamic algorithms require only a linear amount of additional memory.

We also contribute an algorithm to update the group of nodes with the highest group closeness after an edge insertion. It is based on Bergamini and Meyerhenke's improved version of the algorithm by Chen et al. The basic idea is to verify whether the choices of the greedy algorithm are still valid on the modified graph with as little computational expense as possible. Once the greedy algorithm would choose a different node than on the previous graph, the dynamic algorithm discards all information from the previous run and falls back to the static algorithm.

Both dynamic algorithms are based on existing implementations in NetworKit which is an open-source toolkit for high-performance network analysis~\cite{staudt2014networkit}.
%% content.tex
%%

%% ==============
\chapter{Preliminaries}
\label{ch:preliminaries}

In this thesis, $G = (V, E)$ usually denotes a simple, unweighted graph with a set of nodes $V$ and a set of edges $E$ between the nodes. An undirected edge is a subset of $V$ with exactly two distinct elements. A directed edge is an element of $V \times V$. The distance between two nodes $u$ and $v$ is denoted by $d(u, v)$. The set of neighbor nodes for $v$ is denoted by $N(v) := \{u : \exists (v, u) \in E\}$.

We now want to define the concept of closeness centrality. The most simple definition, while only useful for connected graphs, is the following:

\begin{definition}
\label{def:closenessConnected}
Let $G = (V, E)$ be a connected, unweighted graph. The closeness centrality of a node $v \in V$ is defined as
\begin{align*}
	c(v) = \frac{|V| - 1}{\sum_{u \in V}{d(v, u)}}.
\end{align*}
\end{definition}

In the disconnected case, there are node pairs without a path between them. Using Definition~\ref{def:closenessConnected} and assuming $d(u, v) = \infty$ for such node pairs, the sum over all the distances in the denominator would blow up and make the resulting closeness centralities useless. To solve this problem, we first define $R(v) \coloneqq \{u \in V : u \text{ is reachable from }v \text{ in } G\}$ and $r(v) \coloneqq |R(v)|$. This leads to a generalized version of Definition~\ref{def:closenessConnected}:
\begin{align}
	c(v) = \frac{r(v) - 1}{\sum_{u \in R(v)}{d(v, u)}}. \label{eq:closenessIncomplete}
\end{align}
However, this definition does not differentiate between central nodes in small components and central nodes in large components of a graph. Intuitively, a node $v$ with $r(v) = 2000$ and a total distance of $4000$ to all reachable nodes is more central than a node $w$ with $r(w) = 20$ and a total distance of $40$. In order to give preference to nodes in large components, we scale Equation~\ref{eq:closenessIncomplete} by $\frac{r(v)}{|V| - 1}$. 

\begin{definition}
\label{def:closenessGeneral}
Let $G = (V, E)$ be an unweighted graph. The closeness centrality of a node $v \in V$ is defined as
\begin{align*}
	c(v) = \frac{r(v) - 1}{\sum_{u \in R(v)}{d(v, u)}} \cdot \frac{r(v)}{n - 1}.
\end{align*}
\end{definition}

Another approach to handle disconnected graphs is called \emph{harmonic centrality}.

\begin{definition}
\label{def:harmonicCentrality}
Let $G = (V, E)$ be an unweighted graph. The harmonic centrality of a node $v \in V$ is defined as
\begin{align*}
	h(v) = \sum_{u \in V}{\frac{1}{d(v, u)}}.
\end{align*}
\end{definition}
If there is no path between $u$ and $v$, the inverse distance is set to $0$.

In this thesis, we will sometimes use \emph{harmonic centrality} and \emph{closeness centrality} interchangeably.
\chapter{Static Top-k closeness}
\label{ch:staticCloseness}

Computing the exact closeness centrality of each node in a graph requires solving the \emph{all-pair-shortest-path} (APSP) problem . A simple approach is to compute a breadth-first search (or Dijkstra's algorithm) from each node. This results in a time complexity of $\mathcal{O}(|V| \cdot (|V| + |E|))$ for unweighted graphs and $\mathcal{O}(|V| \cdot (|V|\log{|V|} + |E|))$ for weighted graphs. The Floyd-Warshall algorithm has a time complexity of $\mathcal{O}(n^3)$~\cite{Floyd:1962:A9S:367766.368168}; Johnsons's algorithm for weighted graphs without negative cycles has a time complexity of $\mathcal{O}(|V| \cdot (|V|\log{|V|} + |E|))$~\cite{Johnson:1977:EAS:321992.321993}. There are also algorithms that are based on Fast Matrix Multiplication~\cite{zwick2002all,williams2012multiplying} which solve the problem in $\mathcal{O}(n^{2.3727})$. In practice, methods using graph traversal such as breadth-first search or Dijkstra's algorithm are preferred because real-world graphs are often sparse. Fast Matrix Multiplication algorithms also require more memory. 

\section{Efficient algorithms for closeness centrality}
\label{sec:closenessAlgorithms}

There have been some ideas to reduce the effort to compute closeness centralities in large networks.

\subsection{Approximation of closeness centrality}
Eppstein and Wang propose a fast approximation algorithm in~\cite{eppstein2001fast} for large networks exhibiting the \emph{small world phenomenon}. It provides an $(1 + \epsilon)$-approximation in near-linear time. The algorithm works as follows:

\begin{enumerate}
	\item Let $k$ be the number of iterations to obtain the desired error bound
	\item In iteration $i$, pick vertex $v_i$ uniformly at random from $G$ and solve the SSSP problem with $v_i$ as the source
	\item Let
		\[
			\hat{c}_u = \frac{1}{\sum_{i = 1}^{k}{\frac{n \cdot d \cdot (v_i, u)}{k \cdot (n - 1)}}} 	
		\] be the centrality estimator for vertex $u$.
\end{enumerate}
Eppstein and Wang show that the expected value of $\frac{1}{\hat{c}_u}$ is equal to $\frac{1}{c_u}$. Using Hoeffding's bound, they also show that for $k = \mathcal{O}(\frac{\log{n}}{\epsilon^2})$ the additive error is at most $\Delta \epsilon$ with high probability, where $\Delta$ denotes the diameter of the graph. The total runtime of the algorithm is $\mathcal{O}\left(\frac{\log{n}}{\epsilon^2} (n \log{n} + m)\right)$ for weighted graphs. For unweighted graphs it is $\mathcal{O}\left(\frac{\log{n}}{\epsilon^2} \cdot (n + m)\right)$.

\subsection{Top-k closeness centrality}
\label{sec:topKClosenessComplex}
For some real-world applications, it is unnecessary to know the exact closeness centrality and the exact ranking of unimportant nodes. In these cases, it might be enough to compute a list of the $k$ nodes with the highest exact closeness centrality. For all the other nodes it is enough to provide an upper bound that is lower than the exact closeness centrality of the $k$-th most central node.

Borassi et al. propose an efficient algorithm for the problem in \cite{borassi2015fast} which works especially well on complex networks. Angriman presents an adaptation of the algorithms to work with harmonic centrality~\cite{angriman2016efficient}. Bergamini et al. present some modifications for street networks (i.e. graphs with large diameter) in \cite{bergamini2016computing}. Since the dynamic algorithm presented in this thesis is based on these algorithms, we will provide detailed descriptions and analysis. The following explanation is based on the version of the algorithm for harmonic centrality to avoid redundancy and since it is easily applicable to disconnected graphs. However, the adaptation of these algorithms to closeness centrality is fairly straightforward. 

\subsubsection{Upper bounds for the closeness centrality of a node}
\label{sec:borassiUpperBound}
The algorithm by Borassi et al. is briefly outlined in Algorithm~\ref{alg:borassiStatic}. The main idea is to run a BFS from each node in the graph, but abort the search once it is clear that the node does not belong to the $k$ nodes with highest closeness. During the BFS from a specific node $v$, the algorithm keeps track of an upper bound $\widetilde{h}(v)$ for the harmonic closeness of that node.

For that purpose, the algorithm keeps track of the current level $d$ of the search, that is the current distance to the source node $v$ of the search. When the BFS reaches a new level, the upper bound $\widetilde{h}(v)$ is recomputed. If the new upper bound is smaller than the known exact harmonic centrality of the $k$-th node in the list, the search is aborted. 

Let $\Gamma_d$ denote the set of nodes on level $d$ from $v$, $\gamma_d$ the number of nodes on level $d$, i.e. $\gamma_d = |\Gamma_d|$. For each level, the algorithm computes an upper bound $\widetilde{\gamma}_{d+1} = \sum_{u \in \Gamma_d}{deg(u)}$ for the number of nodes on level $d + 1$. Let $r(v)$ denote the number of nodes reachable from $v$. Let $n_d$ denote the number of nodes up to level $d$. Let $h_d(v)$ denote the harmonic closeness based on all nodes up to level $d$ from $v$, i.e. $h_d(v) = \sum_{d(v, u) \leq d}{\frac{1}{d(v, u)}}$.

During the BFS, the algorithms sums up the inverse distances of visited nodes. After visiting all nodes on level $d$, the resulting sum is $h_d(v)$. For the upper bound, we start with
\begin{align}
	h(v) \leq h'(v) = h_d(v) + \frac{\gamma_{d+1}}{d+1} + \frac{r(v) - n_{d+1}}{d+2}.
\end{align}

Basically, all nodes on level $d + 1$ contribute $\frac{1}{d+1}$ to the harmonic centrality. All remaining unvisited nodes have at least distance $d + 2$.

Since $n_{d + 1} = n_d + \gamma_{d+1}$, we can write
\begin{align}
	         h'(v)   &= h_d(v) + \frac{\gamma_{d+1}}{d+1} + \frac{r(v) - n_{d} - \gamma_{d+1}}{d+2}
\end{align}
We can now replace $\gamma_{d+1}$ with its upper bound $\widetilde{\gamma}_{d+1}$, which is computed after visiting all nodes on level $d$.
\begin{align}
	          h'(v)  &\leq h_d(v) + \frac{\widetilde{\gamma}_{d+1}}{d+1} + \frac{r(v) - n_{d} - \widetilde{\gamma}_{d+1}}{d+2} \label{eq:harmonicClosenessEstimate} \\
	                 &= h_d(v) + \frac{(d + 2) \cdot \widetilde{\gamma}_{d+1} + (d+1) \cdot \left(r(v) - n_{d} - \widetilde{\gamma}_{d+1}\right)}{(d+1) \cdot (d+2)} \\
	                 &= h_d(v) + \frac{\widetilde{\gamma}_{d+1} + (d+1) \cdot \left(r(v) - n_{d} \right)}{(d+1) \cdot (d+2)} \\
  \widetilde{h}(v)   &= h_d(v) + \frac{\widetilde{\gamma}_{d+1}}{(d+1) \cdot (d+2)} + \frac{r(v) - n_{d}}{d+2} \label{eq:harmonicClosenessUpperBound}
  \end{align}

\subsubsection{Computing the number of reachable nodes}
\label{sec:reachableNodes}
The number of reachable nodes $r(v)$ in Equation~\ref{eq:harmonicClosenessUpperBound} can be computed in a preprocessing step. For undirected graphs, the number of reachable nodes from $v$ is equal to the size of the connected component containing $v$. Algorithm~\ref{alg:connectedComponents} computes the number of reachable nodes for each node in $\mathcal{O}(n + m)$ with a single BFS.

\paragraph{Directed graphs}
 For directed graphs, $r(v)$ cannot be computes as easily if the graph is not strongly-connected. Instead, an upper bound for $r(v)$ is computed. The strongly-connected components of a graph $G$ can be computed in linear time with Tarjan's algorithm~\cite{tarjan1972depth}. Let $\mathcal{G = (V, E)}$ denote the graph of the strongly-connected components of $G$, with
 
 \begin{itemize}
 	\item $\mathcal{V}$ as the set of strongly-connected components in $G$
 	\item $C \subset V$ and $D \subset V$ denoting strongly-connected components
 	\item $(C, D) \in \mathcal{E} \iff \exists (u, v) \in E, \text{ such that } u \in C \land v \in D$ 
 	\item the weight $w(C) = |C|$ for each strongly-connected component $C$.
 \end{itemize}
 
 Tarjan's algorithm already provides a topologically sorted list of the strongly-connected components of $G$. Components without any outgoing edges, so-called \emph{sinks}, are placed at the end of the list. The list is then processed in reverse order and an upper bound $\omega(C)$ for the number of reachable nodes from that component can be computed with 
 \begin{align*}
 	\omega(C) = w(C) + \sum_{(C, D) \in \mathcal{E}}{\omega(D)}.
 \end{align*}
Since the components are traversed in reverse topological order, the upper bound of each neighbor component is already known. Please note that there are no circles in this graph of strongly-connected components. Otherwise, all components comprising the circle would belong to one larger strongly-connected components.

\begin{algorithm2e}[h!]
 \label{alg:connectedComponents}
 \KwData{$G = (V, E)$}
 \KwResult{$r(v)$ for each node $v$}
 $P(v) \gets$ \texttt{none} for each node $v$ \\
 $i \gets 0$ \\
 
 \ForAll{$v \in V$}{
   	\If{$P(v) = $ \texttt{none}}{
	   Run BFS from $v$ and set $P(u) = i$ for each visited node $u$, keep track of the component sizes \\
	   $i \gets i + 1$
   	}
 }
 
 \ForAll{$v \in V$}{
 	$r(v) \gets componentSizes[P(v)]$
 }
 
 \caption{Computing connected components in an undirected graph}
\end{algorithm2e}

\subsubsection{Pruned breadth-first search}
In order to compute the closeness centrality of a single node in an unweighted graph, a complete breadth-first search starting from that node is necessary. The algorithm only requires the exact closeness centrality of the $k$ most central nodes. Therefore, the breadth-first searches can be aborted early for many nodes, once the upper bound obtained with Equation~\ref{eq:harmonicClosenessUpperBound} is smaller than the exact closeness centrality of the $k$-th most central node among the already processed nodes. Algorithm~\ref{alg:bfsCutStatic} outlines the steps necessary to perform such a pruned BFS. The algorithm is structured like a standard BFS, but also keeps track of the closeness centrality of the source node. The algorithm sums the inverse distances of all visited nodes in $h$ (Line~\ref{alg:inverseDistanceSum}). Once the BFS reaches a new level (Line~\ref{alg:newLevel}), $\widetilde{h}$ is updated (Line~\ref{alg:upperBoundUpdate}). If the new upper bound is smaller than the known exact closeness $x_k$ of the $k$-th node, the algorithm returns with the last computed upper bound and a flag that indicates that the value is not exact. The algorithm also keeps track of an upper bound $\widetilde{\gamma}$ for the number of nodes on level $d + 1$ by summing the out-degrees of all nodes on level $d$ (Line~\ref{alg:gammaUpdateDirected} and Line~\ref{alg:gammaUpdateUndirected}). If the search encounters a node $w$ that has already been marked as visited, that node must be on a level $l < d$. It is then possible to compute a new, lower upper bound $\widetilde{h}$ in some cases (Line~\ref{alg:boundAdjustment}).

\begin{algorithm2e}[h!]
 \label{alg:bfsCutStatic}
 \KwData{$G = (V, E), v, x_k$}
 \KwResult{A tuple $(h, isExact)$ with $isExact = $\texttt{false} if $h$ is only an upper bound for the exact harmonic centrality.}
 Create queue $Q$ \\
 $Q$.enqueue($v$) \\
 Mark $v$ as visisted \\
 $d \gets 0; h \gets 0; \widetilde{\gamma} \gets 0; nd \gets 0$ \\
 
 \While{!Q.isEmpty}{
   $u \gets$ Q.dequeue() \\
   \If{$d(v, u) > d$}{ \label{alg:newLevel}
     $d \gets d + 1$ \\
   	 $r \gets r(u)$ \\
   	 $\widetilde{h} \gets h + \frac{\widetilde{\gamma}}{(d + 1) \cdot (d + 2)} + \frac{r - n_d}{d + 2}$ \label{alg:upperBoundUpdate} \\
   	 \If{$\widetilde{h} \leq x_k$}{
   	   return $(\widetilde{h}, \texttt{false})$
   	 }
   }
   \ForAll{$w \in N(u)$}{
     \uIf{$w$ is not marked as visited}{
       Mark $w$ as visited \\
       $Q$.enqueue($w$) \\
       $n_d \gets n_d + 1$ \\
       $pred[w] \gets u$ \\
       $d(v, w) \gets d(v, u) + 1$ \\
       $h \gets h + \frac{1}{d(v, w)}$ \label{alg:inverseDistanceSum} \\
       \eIf{$G$ is directed}{
       		$\widetilde{\gamma} \gets \widetilde{\gamma} + outdegree(w) - 1$\label{alg:gammaUpdateDirected}
       }{
       		$\widetilde{\gamma} \gets \widetilde{\gamma} + outdegree(w)$\label{alg:gammaUpdateUndirected}
       }
        
     }
     \uElseIf{$d(v, w) > 1 \land pred[u] \neq w$}{ \label{alg:boundAdjustment}
     	$\widetilde{h} \gets \widetilde{h} - \frac{1}{d + 1} + \frac{1}{d + 2}$ \\
     	\If{$\widetilde{h} \leq x_k$}{
     		return $(\widetilde{h}, \texttt{false})$
     	}
     }
   }
   return $(h, \texttt{true})$
 }
 
 \caption{BFScut()}
\end{algorithm2e}

\subsubsection{Algorithm outline}
Algorithm~\ref{alg:borassiStatic} is used to compute the list of the $k$ nodes with the highest closeness in $G$. The preprocessing in Line~\ref{alg:borassiStaticPreprocessing} is used to compute the number of reachable nodes for each node in the graph in linear time (see Section~\ref{sec:reachableNodes}). The algorithm then iterates over all nodes in decreasing order of degree (Line~\ref{alg:borassiLoop}) and starts a pruned BFS from each node. If the pruned BFS returns an exact closeness value, and if the value is larger than the current $k$-th largest value, it can be added to the list with the most central nodes. $x_k$, representing the $k$-th largest value, can then also be updated.


\begin{algorithm2e}[h!]
 \label{alg:borassiStatic}
 \KwData{$G = (V, E)$}
 \KwResult{A list with the $k$ nodes with the highest closeness centrality}
 Compute $r(v) \forall v \in V$ \label{alg:borassiStaticPreprocessing}\\
 $\texttt{Top} \gets $ empty priority queue
 $h(v) \gets 0$ for each node $v$\\
 $isExact(v) \gets $ \texttt{false} for each node $v$\\
 $x_k \gets 0$\\

 \ForAll{$v \in V$ in decreasing order of degree}{ \label{alg:borassiLoop}
  	$(h, isExact) \gets $ \texttt{BFSCut}($v, x_k$) \\
  	$h(v) \gets h$\\
  	$isExact(v) \gets isExact$\\ 
  	\If{$isExact \land h > x_k$}{
  	   \texttt{Top.insert(h, v)} \\
  	   \If{$\texttt{Top.size()} > k$}{
	      \texttt{Top.removeMin()}
	   }
	   \If{$\texttt{Top.size()} = k$}{
	     $x_k \gets \texttt{Top.getMin()}$
	   }
	}
 }
 \caption{Static computation of the $k$ nodes with the highest closeness}
\end{algorithm2e}

\subsubsection{Networks with large diameter}
\label{sec:largeDiameterStatic}
The previously described algorithm works well for complex (social) networks with small diameter. Bergamini et al. present a different approach to solve the problem faster on networks with large diameter, for instance street networks~\cite{bergamini2016computing}. The basic idea is to always run a complete breadth-first search from a source node and then use the number of nodes on each level to compute upper bounds for the closeness centrality of each other node in the graph. The nodes are kept in a list sorted in decreasing order by the corresponding upper bound for their closeness centrality. We provide a detailed explanation of this strategy in the following paragraphs. 

\paragraph{Level-based upper bounds}
It is possible to compute upper bounds for the closeness centrality of all nodes in a graph with a single BFS. Let $G$ denote an undirected graph and $s$ the source node of the BFS. A node $v$ is on level $i$ ($l(v) = i$) if $d(s, v) = i$ and we write $v \in \Gamma_i(s) \iff d(s, v) = i$ (see Section~\ref{sec:borassiUpperBound}). The distance between two arbitrary nodes $v \in \Gamma_i(s)$ and $w \in \Gamma_j(s)$ for $i \leq j$ is at least $j - i$. If $d(v, w)$ was smaller than $j - i$, $w$ would have been discovered earlier and its level would be $i + d(v, w) < j$. This is a contradiction to the assumption that $w$ is on level $j$. With this discovery, it is possible to obtain an upper bound for the harmonic closeness of each node $v \in V$:
\begin{align*}
	h(v) \leq \widetilde{h}(v) = \sum_{w \in R(s)}{\left| \frac{1}{d(s, w)} - \frac{1}{d(s, v)} \right| }.
\end{align*}
This bound can be improved further. The degree of a node $v$ is equal to the number of nodes $w$ at distance 1. All the other nodes with $\left| \frac{1}{d(s, w)} - \frac{1}{d(s, v)}\right| \leq 1$ must have at least distance 2. This leads to
\begin{align}
	\widetilde{h}(v) ={} &1\cdot deg(v) \nonumber \\
	               &+ \frac{1}{2} \cdot \left( \left|\left\{ w \in R(s) : \left| d(s, w) - d(s, v) \right| \leq 1 \right\} \right| - deg(v) - 1 \right) \nonumber \\
	               &+ \sum_{\substack{w \in R(s) \\ |d(s, w) - d(s, v)| > 1}}{|\frac{1}{d(s, w)} - \frac{1}{d(s, v)}|}. \label{eq:levelBasedUpperBound} 
\end{align}
After the BFS, the number of nodes $\gamma_j$ on each level $j$ is known. With that information, Equation~\ref{eq:levelBasedUpperBound} can be rewritten to
\begin{align}
	\widetilde{h}(v) = \frac{1}{2} \cdot \left(\sum_{|j - d(s, v)| \leq 1}{\gamma_j}\right) + \left(\sum_{|j - d(s, v)| > 1}{\gamma_j \cdot \left|\frac{1}{j - d(s, v)}\right|}\right) - \frac{1}{2} + \frac{1}{2} \cdot deg(v).
\end{align}
For all nodes $v$ with the same distance from $s$, the first three terms of the equation are the same. Therefore, a preliminary \emph{level bound} can be computed once for each level. For each individual node $v$, only $\frac{1}{2}\cdot deg(v)$ needs to be added to get $\widetilde{h}(v)$.

For directed graphs, Equation~\ref{eq:levelBasedUpperBound} does not hold. Consider two nodes $v$ and $w$ with $l(w) < l(v)$. In the undirected case, it is possible to infer a lower bound for the distance between the two nodes. This is not possible in the directed case because there could be a shortcut between $v$ and $w$ as shown in Figure~\ref{fig:directedBFSbound}. For all nodes $w$ with $l(w) < l(v)$ and $w \notin N(v)$, we can only assume that the distance must be at least $2$. This leads to slightly less tight upper bounds for directed graphs. Modifying Equation~\ref{eq:levelBasedUpperBound} accordingly leads to
\begin{align}
     \widetilde{h}_{directed}(v) ={} &\frac{1}{2} \cdot deg(v) \nonumber \\
      &+ \frac{1}{2} \cdot \left( \left|\left\{ w \in R(s) : d(s, w) - d(s, v) \leq 1 \right\} \right| \right) \nonumber \\
      &+ \sum_{\substack{w \in R(s) \\ d(s, w) - d(s, v) > 1}}{\frac{1}{d(s, w)} - \frac{1}{d(s, v)}}. \label{eq:levelBasedUpperBoundDirected} 
\end{align}



\begin{figure}[h!]
\centering
\begin{tikzpicture}
	\tikzset{edge/.style = {->,> = latex'}}
	
    \node[main node,fill=white] (1) {$s$};
    \node[main node] (2) [below left = 1.3cm and 0.8cm of 1]  {$w$};
    \node[main node,fill=white] (3) [below right = 1.3cm and 0.8cm of 1] {};
    \node[main node,fill=white] (4) [below left = 1.3cm and 0.8cm of 2] {};
    \node[main node,fill=white] (5) [below left = 1.3cm and 0.8cm of 4] {};
    \node[main node] (6) [below left = 1.3cm and 0.8cm of 5] {$v$};
    
    
	\draw[edge] (1) to (2);
	\draw[edge] (1) to (3);
	\draw[edge] (2) to (4);
	\draw[edge] (4) to (5);
	\draw[edge] (5) to (6);
	\draw[edge] (6) to[bend right] (3);
	\draw[edge] (3) to (2);
	
\end{tikzpicture}
\caption{BFS tree}{Distances in directed graphs are not symmetric. The distance between $w$ and $v$ is $3$ in the BFS tree rooted in $s$, while the distance in the actual directed graph is $2$.}
\label{fig:directedBFSbound}
\end{figure}

\paragraph{Algorithm outline}
The algorithm to compute the $k$ nodes with highest closeness in networks with large diameters is outlined in Algorithm~\ref{alg:streetNetworksStatic}. The algorithm uses a priority queue \texttt{Q} which contains all unprocessed nodes, sorted by decreasing priority. It is required that the priority queue makes it possible to modify the priorities of existing elements efficiently. Initially, the degree of each node can be used for its priority. Alternatively, Bergamini et al. propose a preprocessing algorithm which computes an initial upper bound for the closeness centrality of each node based on its extended neighborhood, that is, all nodes up to a configurable distance $l$. The algorithm maintains an array \texttt{score} which contains the closeness centralities of each node. Each entry is initialized to $\infty$ (this means \texttt{std::numeric\_limits<double>::max()} in a C++ implementation). The priority queue \texttt{Top} maintains the $k$ nodes with highest closeness in increasing order.

The algorithm processes the nodes in \texttt{Q} in decreasing order of priority (Line~\ref{alg:streetNetworksLoop}). In each iteration, the node $v$ with the maximum priority is extracted from \texttt{Q} (Line~\ref{alg:streetNetworksExtraction}). We call $v$ the \emph{active} node. If the known upper bound for the closeness of the node is lower than the exact known value for the $k$-th most central node, the algorithm can be aborted and the list with the $k$ nodes is returned. Otherwise, the \texttt{updateBounds} function is called with $v$ as the source node (Line~\ref{alg:updateBoundsCall}).

The \texttt{updateBounds} function is outlined in Algorithm~\ref{alg:updateBounds}. It computes a complete BFS from the supplied source node $s$ and counts the number of nodes $\gamma_i$ on each level $i$. The complete BFS from $s$ makes it possible to compute the exact closeness of $s$. The resulting upper bound for each other node in the graph is computed using Equation~\ref{eq:levelBasedUpperBound} in the undirected case and Equation~\ref{eq:levelBasedUpperBoundDirected} in the directed case.

After the call to \texttt{updateBounds} returns, Algorithm~\ref{alg:streetNetworksStatic} updates \texttt{score}. If the upper bound for a node $u$ returned by the function is smaller than the previous upper bound stored in \texttt{score[$u$]}, the array is updated with the new, tighter upper bound (Line~\ref{alg:updateScores}). In order to keep the list of the $k$ most central up-to-date, the active node $v$ is inserted with its exact closeness centrality as the priority into \texttt{Top} if the exact closeness centrality of $v$ is larger than that of the $k$-th node in \texttt{Top} (Line~\ref{alg:updateTop}). If \texttt{Top} contains more than $k$ elements, the node with the smallest closeness is removed from the list.

\todo{Mention parallelization}

\begin{algorithm2e}[h!]
 \label{alg:streetNetworksStatic}
 \KwData{$G = (V, E)$}
 \KwResult{A list with the $k$ nodes with the highest closeness}
 \texttt{Preprocessing}($G$) \\ \label{alg:streetNetworksPreprocessing}
 \texttt{Q} $\gets V$, sorted by decreasing degree or a precomputed upper bound \\
 \texttt{Top} $\gets []$ \\
 \texttt{score} $\gets $ array indexed by node ID storing the current upper bounds for all nodes 
 
 \For{$v \in V$}{\texttt{score[$v$]} $\gets \infty$}
 
 \While{$Q$ is not empty}{ \label{alg:streetNetworksLoop}
   $v \gets \texttt{Q.extractMax()}$ \label{alg:streetNetworksExtraction} \\
   
   \If{$|\texttt{Top}| \geq k \land $\texttt{score[$v$]}$ \leq $ \texttt{Top[$k$]}}{
     return \texttt{Top}
   } 
   
   \texttt{levelBasedBounds} $\gets $ \texttt{updateBounds($v$)} \label{alg:updateBoundsCall} \\
   $\texttt{score[v]} \gets \texttt{levelBasedBounds[v]}$ \\
   
   \ForAll{$w \in V$}{ \label{alg:updateScores}
     \If{\texttt{levelBasedBounds}[$w$] $<$ \texttt{score[$w$]}}{
       \texttt{score[$w$]} $\gets $ \texttt{levelBasedBounds}[$w$]
     }
   }
   
   \If{\texttt{score[$v$]} > \texttt{Top[$k$]}}{ \label{alg:updateTop}
      add $v$ to \texttt{Top} \\
   	  sort \texttt{Top} by \texttt{score} and reduce it to at most $k$ elements
   }
   re-order \texttt{Q} according to the new values in \texttt{score}
    
 }
  
 \caption{Static computation of the $k$ nodes with the highest closeness centrality in networks with large diameter}
\end{algorithm2e}

\begin{algorithm2e}[h!]
	\label{alg:updateBounds}
 \KwData{$G = (V, E), s \in V$}
 \KwResult{The exact closeness centrality of $s$, upper bounds for the closeness of all other nodes}
  
  $d \gets \texttt{BFSfrom}(s)$ \\ \label{alg:bfsFromSource}
  $d_{max} \gets max_{v \in V}d(s, v)$ \\

  $\forall i \in [1..d_{max}]: \gamma_i \gets |\{v : d(s, v) = i\}|$ \\
  \ForAll{$v \in V$}{
  	$S[v] \gets \frac{1}{2} \cdot \left(\sum_{|j - d(s, v)| \leq 1}{\gamma_j}\right) + \left(\sum_{|j - d(s, v)| > 1}{\gamma_j \cdot \left|\frac{1}{j - d(s, v)}\right|}\right) - \frac{1}{2} + \frac{1}{2} \cdot deg(v)$
  }
  $S[s] \gets \sum_{u \in R(s)}{\frac{1}{d(s, u)}}$ \\
  
  return $S$
  
 \caption{The \texttt{updateBounds} function computes the exact closeness centrality of the supplied source node $s$ and provides upper bounds for the closeness centrality of all other nodes in the graph.}
\end{algorithm2e}






%% conclusion.tex
%%

%% ==================
\chapter{Conclusion}
\label{ch:conclusion}
%% ==================

We have developed and implemented dynamic algorithms for Top-$k$ closeness centrality for both complex networks and street networks. We initially limit the recomputation of closeness centralities or upper bounds thereof to nodes actually affected by an edge modification. Additionally, we can exploit the fact that we only need to compute the exact closeness centralities for the $k$ most central nodes. We can keep the upper bounds of nodes which are far away from an edge insertion in place and can update the upper bounds of so-called boundary nodes cheaply. At last, we can compute the maximum improvement of the closeness centrality for each node. If the resulting upper bound is smaller than the $k$-th largest closeness centrality, we can use the improvement bound and do not need a new breadth-first search.

Our dynamic algorithm is between 35 and 45 times faster on average than the static algorithm for edge insertions on undirected complex networks. On directed complex networks, the dynamic algorithm is between 15 and 20 times faster. In the best case, our dynamic algorithm is up to 700 times faster for some edges.

We achieve even higher speedups on street networks, improving by a factor of more than 100 on average with directed street networks. In some cases, our dynamic algorithm is up to three orders of magnitudes faster than the static algorithm. Our analysis shows that the speedup is the result of a much lower number of complete breadth-first searches compared to the static algorithm.

We have also implemented an incremental algorithm based on the greedy algorithm for the Maximum Closeness Centrality Group Identification problem. We make use of our dynamic algorithm for Top-$k$ closeness centrality and try to replicate the results of the static algorithm with as little computational effort as possible. We avoid recomputing the marginal gains for as many nodes as possible. Once any node in the group has changed, we fall back to the static algorithm.

The dynamic algorithm is between 8 and 15 times faster than the static algorithm on average, and up to 800 times faster for some inserted edges on the larger networks we have tested.

\paragraph{Future work}
We have already implemented a simple optimization to speed up the computation of the number of reachable nodes for each node in undirected graphs after an edge insertion. As we have seen in our experimental evaluation, recomputing an upper bound for the number of reachable nodes in directed graphs after edge modifications is a significant bottleneck. There are algorithms that maintain reachability information in directed graphs after edge removals~\cite{lkacki2013improved,chechik2016decremental}. Since these algorithms maintain exact reachability information for single source nodes, it might be impractical to use them for our purposes. However, it might be practical to maintain upper bounds for the number of reachable nodes.

Our dynamic algorithm is currently designed to update the ranking of the $k$ most central nodes after every single edge modification. However, it might not be necessary to update the Top-$k$ nodes after each edge modification. Then it might be possible to coalesce multiple edge insertions or edge removals into a single update operation and still be faster than running the static algorithm. This is called \emph{batch update}. One could compute the set of affected nodes individually for each edge in the batch and compute the union of all these sets. Even far-away nodes could still be skipped for edge insertions if their minimum distance to any of the edges is larger than their previous cutoff level.

In our dynamic algorithm for group closeness, we perform a pruned BFS for each affected node in each iteration of the algorithm. It might be possible to adapt the idea of the level-based improvement bounds for the group closeness algorithm. If it is possible to compute an upper bound for the increase of the marginal gain for all nodes, we could skip the expensive pruned BFS for nodes where we are sure that the new marginal gain will be smaller than the currently known maximum.



%% declaration.tex
%%

%% ==================
\iflanguage{english}
{\chapter{Declaration}}
{\chapter{Erklärung}}
\label{ch:Declaration}
%% ==================

Ich versichere hiermit wahrheitsgemäß, die Arbeit selbstständig verfasst und
keine anderen als die angegebenen Quellen und Hilfsmittel benutzt, die wörtlich
oder inhaltlich übernommenen Stellen als solche kenntlich gemacht und die
Satzung des Karlsruher Instituts für Technologie (KIT) zur Sicherung guter
wissenschaftlicher Praxis in der jeweils gültigen Fassung beachtet zu
haben.

\vspace*{1cm}
\hspace*{4cm} Karlsruhe, den \submissiontime \hspace*{0.5cm}\hrulefill \\
\hspace*{10.5cm} \myname \\




%% ----------------
%% |   Appendix   |
%% ----------------
% IM Style: No additional blank page
% \cleardoublepage

%%% appendix.tex
%%

%% ==============================
%\chapter{Appendix}
%\label{ch:Appendix}
%% ==============================

\appendix

\iflanguage{english}
{\addchap{Appendix}}	% english style
{\addchap{Anhang}}	% german style


\section{First Appendix Section}
		\label{Anhang-Implementierung}
		
\setcounter{figure}{0}

\dots






%% --------------------
%% |   Bibliography   |
%% --------------------
\cleardoublepage
\phantomsection
\addcontentsline{toc}{chapter}{\bibname}

% IM Style
\iflanguage{english}
{\bibliographystyle{plain}}	% english style
{\bibliographystyle{plain}}	% german style

%

% Informatik-Style
%\iflanguage{english}
%{\bibliographystyle{IEEEtranSA}}	% english style
%{\bibliographystyle{babalpha-fl}}	% german style
												  
% Use IEEEtran for numeric references
%\bibliographystyle{IEEEtranSA})

\bibliography{thesis}


\end{document}
